% 第12章图（12.4.8）：Golub--Kahan双对角化的交替递推与矩阵关系 A V_k = U_{k+1} B_k，取 k=3。
% 递推：beta_1 u_1 = b；alpha_i v_i = A^T u_i - beta_i v_{i-1}；beta_{i+1} u_{i+1} = A v_i - alpha_i u_i。
% B_k 的第 i 列为 (alpha_i, beta_{i+1})，位于第 i、i+1 行。
\begin{tikzpicture}[x=1mm,y=1mm,>=Stealth,
  vec/.style={circle,draw=FigureStroke,fill=white,line width=.7pt,minimum size=7mm,inner sep=0pt},
  uvec/.style={vec,draw=FigureBlue,fill=FigureBlueFill},
  vvec/.style={vec,draw=FigureGreen,fill=FigureGreenFill},
  note/.style={font=\FigureNoteFont\fontsize{7.5}{9}\selectfont,text=FigureMuted},
  lab/.style={inner sep=1.2pt},
  box/.style={rectangle,rounded corners=1.8pt,draw=FigureStroke,fill=white,line width=.7pt,inner sep=0pt},
  cell/.style={rectangle,draw=FigureGrid,line width=.4pt,minimum width=5.6mm,minimum height=4.6mm,inner sep=0pt}]
% Preserve the previous canvas while relocating panel labels.
\path[use as bounding box] (-17.833135,-67.672823) rectangle (90.122136,11.562808);
\begin{scope}[yshift=5.5mm]
% ---------- (a) 交替递推 ----------
\node[anchor=north,inner sep=0pt,font=\fontsize{8}{10}\selectfont] at (36,-21.8) {(a)};
\node[note,anchor=east] at (3,0) {数据空间$\mathbb R^m$};
\node[note,anchor=east] at (3,-15) {未知量空间$\mathbb R^d$};
\node[vec] (b) at (7,0) {$\bm b$};
\foreach \i/\x in {1/23,2/45,3/67}{\node[uvec] (u\i) at (\x,0) {$\bm u_{\i}$};}
\foreach \i/\x in {1/34,2/56,3/78}{\node[vvec] (v\i) at (\x,-15) {$\bm v_{\i}$};}
\draw[->,FigureStroke,line width=.8pt] (b)--node[lab,above]{$\beta_1$}(u1);
\foreach \i in {1,2,3}{
  \draw[->,FigureGreen,line width=1pt] (u\i)--node[lab,left=.6mm,text=FigureGreen]{$\alpha_{\i}$}(v\i);}
\foreach \i [evaluate=\i as \j using int(\i+1)] in {1,2}{
  \draw[->,FigureBlue,line width=1pt] (v\i)--node[lab,right=.6mm,text=FigureBlue]{$\beta_{\j}$}(u\j);}
\draw[->,FigureBlue,line width=1pt] (v3)--++(8,8) node[lab,right,text=FigureBlue]{$\beta_4$};
% ---------- (b) 矩阵关系 ----------
\begin{scope}[yshift=-37mm]
\node[anchor=north,inner sep=0pt,font=\fontsize{8}{10}\selectfont] at (20,-33) {(b)};
\def\ytop{-6}
\draw[box] (-12,\ytop) rectangle ++(13,-19); \node at (-5.5,\ytop-9.5) {$\bm A$};
\draw[box,draw=FigureGreen,fill=FigureGreenFill] (3.5,\ytop) rectangle ++(8,-13); \node at (7.5,\ytop-6.5) {$\bm V_3$};
\node at (15.5,\ytop-9.5) {$=$};
\draw[box,draw=FigureBlue,fill=FigureBlueFill] (20,\ytop) rectangle ++(11,-19); \node at (25.5,\ytop-9.5) {$\bm U_4$};
\node[note] at (-5.5,\ytop-22.5) {$m\times d$};
\node[note] at (7.5,\ytop-16.5) {$d\times3$};
\node[note] at (25.5,\ytop-22.5) {$m\times4$};
\foreach \r in {1,...,4}{\foreach \c in {1,2,3}{
  \node[cell] (B\r\c) at ({35+(\c-.5)*5.6},{\ytop-(\r-.5)*4.6}) {};}}
\foreach \i in {1,2,3}{
  \node[cell,draw=FigureGreen,fill=FigureGreenFill] at (B\i\i) {$\alpha_{\i}$};}
\foreach \i [evaluate=\i as \j using int(\i+1)] in {1,2,3}{
  \node[cell,draw=FigureBlue,fill=FigureBlueFill] at (B\j\i) {$\beta_{\j}$};}
\draw[FigureStroke,line width=.7pt] ($(B11.north west)+(.8,0)$)--(B11.north west)--(B41.south west)--++(.8,0);
\draw[FigureStroke,line width=.7pt] ($(B13.north east)+(-.8,0)$)--(B13.north east)--(B43.south east)--++(-.8,0);
\node[note] at ($(B42.south)+(0,-3.2)$) {$\bm B_3$：$4\times3$};
\end{scope}
\end{scope}
\end{tikzpicture}
